% Options for packages loaded elsewhere
\PassOptionsToPackage{unicode}{hyperref}
\PassOptionsToPackage{hyphens}{url}
\documentclass[
]{article}
\usepackage{xcolor}
\usepackage{amsmath,amssymb}
\setcounter{secnumdepth}{-\maxdimen} % remove section numbering
\usepackage{iftex}
\ifPDFTeX
  \usepackage[T1]{fontenc}
  \usepackage[utf8]{inputenc}
  \usepackage{textcomp} % provide euro and other symbols
\else % if luatex or xetex
  \usepackage{unicode-math} % this also loads fontspec
  \defaultfontfeatures{Scale=MatchLowercase}
  \defaultfontfeatures[\rmfamily]{Ligatures=TeX,Scale=1}
\fi
\usepackage{lmodern}
\ifPDFTeX\else
  % xetex/luatex font selection
\fi
% Use upquote if available, for straight quotes in verbatim environments
\IfFileExists{upquote.sty}{\usepackage{upquote}}{}
\IfFileExists{microtype.sty}{% use microtype if available
  \usepackage[]{microtype}
  \UseMicrotypeSet[protrusion]{basicmath} % disable protrusion for tt fonts
}{}
\makeatletter
\@ifundefined{KOMAClassName}{% if non-KOMA class
  \IfFileExists{parskip.sty}{%
    \usepackage{parskip}
  }{% else
    \setlength{\parindent}{0pt}
    \setlength{\parskip}{6pt plus 2pt minus 1pt}}
}{% if KOMA class
  \KOMAoptions{parskip=half}}
\makeatother
\usepackage{longtable,booktabs,array}
\newcounter{none} % for unnumbered tables
\usepackage{calc} % for calculating minipage widths
% Correct order of tables after \paragraph or \subparagraph
\usepackage{etoolbox}
\makeatletter
\patchcmd\longtable{\par}{\if@noskipsec\mbox{}\fi\par}{}{}
\makeatother
% Allow footnotes in longtable head/foot
\IfFileExists{footnotehyper.sty}{\usepackage{footnotehyper}}{\usepackage{footnote}}
\makesavenoteenv{longtable}
\setlength{\emergencystretch}{3em} % prevent overfull lines
\providecommand{\tightlist}{%
  \setlength{\itemsep}{0pt}\setlength{\parskip}{0pt}}
\usepackage{bookmark}
\IfFileExists{xurl.sty}{\usepackage{xurl}}{} % add URL line breaks if available
\urlstyle{same}
\hypersetup{
  pdftitle={Where Representations Diverge: Robust Evidence on Modality{,} Dataset{,} and Depth in Cross-Modal Geometry},
  pdfauthor={Abinash Karki (Independent Research)},
  hidelinks,
  pdfcreator={LaTeX via pandoc}}

\title{Where Representations Diverge: Robust Evidence on Modality,
Dataset, and Depth in Cross-Modal Geometry}
\author{Abinash Karki (Independent Research)}
\date{March 2026}

\begin{document}
\maketitle

\begin{center}\rule{0.5\linewidth}{0.5pt}\end{center}

\subsection{Abstract}\label{abstract}

The Platonic Representation Hypothesis proposes that sufficiently
capable models trained on different modalities converge toward a shared
representational geometry. We test that claim using a fully local
22-model experiment on Apple Silicon (M1, 16GB RAM): 8 language models,
10 self-supervised vision models, and 4 vision-language models. The
final evaluation uses 28 concepts (20 base + 8 compounds), 30 images per
concept, image-level caching, image-bootstrap confidence intervals,
Mantel permutation tests with Benjamini-Hochberg FDR correction, source
holdout analysis, prompt-sensitivity controls, and an aligned-layer
protocol. Across 231 model pairs, within-family agreement is strong for
vision-language pairs (median Spearman RSA \(\rho = 0.702\)),
vision-vision pairs (\(\rho = 0.682\)), and language-language pairs
(\(\rho = 0.661\)), but much weaker for language-vision pairs
(\(\rho = 0.259\)) and language-vision-language pairs
(\(\rho = 0.230\)). Only 35/80 language-vision pairs and 5/32
language-vision-language pairs survive FDR correction, versus 43/45
vision-vision pairs and 40/40 vision-vision-language pairs. Robustness
controls materially affect interpretation: leave-one-source-out removal
of ImageNet shifts pairwise RSA by a mean of -0.1564 (max
\(|\Delta| = 0.7080\)), language prompt choice changes cross-modal RSA
by a mean maximum absolute delta of 0.1960, and aligned-layer analysis
shows monotonic growth in mean pairwise RSA from 0.3424 at early depth
to 0.4639 at terminal depth. Yet terminal aligned and baseline results
are nearly identical (mean delta +0.0006), indicating that the depth
effect is about when convergence appears, not whether the final layer
changes the conclusion. The strongest supported claim is therefore not
universal convergence, but modality-asymmetric and dataset-conditional
geometry: vision-language encoders remain structurally closer to vision
than to language, and cross-modal geometry is sensitive to image source,
prompt framing, and representational depth.

\textbf{Keywords}: Platonic Representation Hypothesis, representational
similarity analysis, cross-modal convergence, multimodal models, source
holdout, prompt sensitivity, layer alignment

\begin{center}\rule{0.5\linewidth}{0.5pt}\end{center}

\subsection{1. Introduction}\label{introduction}

\subsubsection{1.1 The Question Behind the
Hypothesis}\label{the-question-behind-the-hypothesis}

When distributional embeddings were found to encode relational
regularities such as ``king - man + woman = queen'' (Mikolov et al.,
2013), they suggested that representation may discover geometry rather
than merely store labels. The Platonic Representation Hypothesis (Huh et
al., 2024) pushes that intuition further: if different model families
are all trained on rich evidence about the same world, then their
internal spaces should converge toward the same underlying structure.

That is a strong claim. It says representational geometry is not
primarily a byproduct of modality, architecture, or objective design,
but a consequence of modeling reality itself. Under that view, a pure
language model and a pure vision model should increasingly agree on how
concepts relate, even if they were never trained together.

\subsubsection{1.2 Why This Matters Beyond
AI}\label{why-this-matters-beyond-ai}

The question matters because it reframes convergence as evidence about
inevitability. If very different systems repeatedly discover similar
structures, then some representational solutions may be mathematically
favored rather than historically accidental. That matters not only for
AI, but for broader questions about perception, abstraction, and
biological constraint.

A useful analogy is convergent evolution. Foveae, sparse coding, and
selective attention recur in biological systems because they solve
recurring problems under recurring constraints. The Platonic hypothesis
asks whether an analogous statement holds for neural representation: do
distinct model families discover the same geometry because the world
itself imposes it?

\subsubsection{1.3 From Positive Results to Boundary
Conditions}\label{from-positive-results-to-boundary-conditions}

The hypothesis is easy to overstate. Apparent convergence can arise from
multiple weaker causes: shared internet distributions, architectural
bias, language-alignment objectives, prompt choices, or evaluation
artifacts. A useful experiment therefore does not only ask whether
convergence exists. It asks what survives after the obvious confounds
are stressed.

This paper is organized around that stricter standard. We retain the
original V1 question, but we answer it using the final 22-model
replication package rather than the earlier prototype runs. The goal is
not to extract the most optimistic reading of the data. It is to
identify the most defensible one.

\subsubsection{1.4 This Paper}\label{this-paper}

We evaluate 22 models across three families:

\begin{itemize}
\tightlist
\item
  8 language models
\item
  10 self-supervised vision models
\item
  4 vision-language image encoders
\end{itemize}

The final dataset contains 20 base concepts, 8 compound concepts, and 30
curated images per concept. We measure cross-model agreement with
Representational Similarity Analysis (RSA) over the 20 base concepts and
supplement it with image-bootstrap confidence intervals, Mantel
permutation tests, source holdout analysis, prompt sensitivity, and an
aligned-layer depth profile.

Our research questions are:

\begin{itemize}
\tightlist
\item
  \textbf{RQ1}: How much representational geometry is shared within and
  across modality families?
\item
  \textbf{RQ2}: Do vision-language encoders bridge language and vision,
  or do they remain predominantly vision-like?
\item
  \textbf{RQ3}: How stable are the conclusions under changes in image
  source, text prompt, and layer depth?
\item
  \textbf{RQ4}: Does compound-concept behavior add independent evidence
  for universal convergence?
\end{itemize}

Our central claim is narrow: the final replication supports robust
within-family convergence and strong vision-to-vision-language coupling,
but it does not support a strong version of universal cross-modal
convergence at the scales tested here.

\begin{center}\rule{0.5\linewidth}{0.5pt}\end{center}

\subsection{2. Related Work}\label{related-work}

Huh et al.~(2024) introduced the Platonic Representation Hypothesis and
argued that stronger models show increasing representational agreement
across language and vision. Their framing is useful because it turns a
philosophical intuition into a measurable claim about geometry.

Representational Similarity Analysis (RSA) provides the relevant tool
(Kriegeskorte et al., 2008). RSA compares the geometry of pairwise
relationships inside two representational spaces rather than the
coordinates of the embeddings themselves, making it appropriate for
comparing models with different dimensionalities and architectures.

The present work differs from frontier-scale Platonic evaluations in
three ways. First, it is deliberately constrained to a local compute
regime, forcing the question toward minimum viable scale rather than
massive capacity. Second, it emphasizes robustness controls that can
materially change interpretation: image bootstrap, source holdout,
prompt sensitivity, and aligned-layer analysis. Third, it uses a broad
family comparison rather than a single bridge-model success case.

This produces a more conservative but more interpretable result: not
simply whether some cross-modal signal exists, but how that signal
compares to within-family agreement and how stable it remains once the
obvious nuisance variables are perturbed.

\begin{center}\rule{0.5\linewidth}{0.5pt}\end{center}

\subsection{3. Methods}\label{methods}

\subsubsection{3.1 Model Selection}\label{model-selection}

The final core roster contains 22 models chosen to remain within the
lab's hardware constraint while spanning distinct modality families and
training objectives.

\textbf{Table 1: Final model roster.}

{\def\LTcaptype{none} % do not increment counter
\begin{longtable}[]{@{}
  >{\raggedright\arraybackslash}p{(\linewidth - 2\tabcolsep) * \real{0.4667}}
  >{\raggedright\arraybackslash}p{(\linewidth - 2\tabcolsep) * \real{0.5333}}@{}}
\toprule\noalign{}
\begin{minipage}[b]{\linewidth}\raggedright
Family
\end{minipage} & \begin{minipage}[b]{\linewidth}\raggedright
Models
\end{minipage} \\
\midrule\noalign{}
\endhead
\bottomrule\noalign{}
\endlastfoot
Language (8) & Falcon3-1B-Instruct-8bit, Granite-3.3-2B-Instruct-8bit,
LFM2-2.6B-Exp-8bit, Qwen2.5-1.5B-Instruct-8bit, Qwen3-0.6B-MLX-8bit,
Qwen3-1.7B-MLX-8bit, Qwen3-4B-MLX-8bit, SmolLM3-3B-8bit \\
Vision SSL (10) & BEiT-base, ConvNeXt-v2, DINOv2-base, DINOv2-small,
DINOv3-ConvNeXt-tiny, Hiera-base, I-JEPA, ViT-MAE-base, ViT-MSN-base,
data2vec-vision \\
Vision-Language (4) & CLIP-ViT-B32, MetaCLIP-B32-400m, SigLIP,
SigLIP2 \\
\end{longtable}
}

The language models cover roughly 0.6B to 4B parameters and include both
MLX and Transformer backends. The vision group spans multiple
self-supervised paradigms, including masked prediction,
self-distillation, joint-embedding prediction, and ConvNeXt-based
encoders. The vision-language group provides the key bridge condition:
models trained with explicit image-language alignment but evaluated here
using their image encoders alone.

\subsubsection{3.2 Concept Set and Image
Construction}\label{concept-set-and-image-construction}

The evaluation uses 28 concepts:

\begin{itemize}
\tightlist
\item
  20 base concepts for RSA
\item
  8 compound concepts for exploratory compositionality analysis
\end{itemize}

The 20 base concepts are:

\begin{itemize}
\tightlist
\item
  cat, dog, bird, fish, elephant
\item
  car, bridge, airplane, mountain, building
\item
  fire, water, forest, ocean, road
\item
  city, space, sun, moon, storm
\end{itemize}

The 8 compound concepts are:

\begin{itemize}
\tightlist
\item
  forest fire, space city, water city, city forest
\item
  mountain road, ocean bridge, city bridge, mountain forest
\end{itemize}

Each concept has exactly 30 images. The final manifest contains full
CLIP-score coverage for all 28 concepts and allocates source coverage at
the concept level across ImageNet (10 concepts), OpenImages (6
concepts), and Unsplash (12 concepts). This source diversity matters
empirically, because later holdout analysis shows that image source is
not a cosmetic detail but a major driver of measured geometry.

\subsubsection{3.3 Embedding Extraction and Layer
Protocol}\label{embedding-extraction-and-layer-protocol}

For language models, concept embeddings are extracted from a short
textual prompt. The canonical template is:

\texttt{The\ concept\ of\ \{concept\}}

Prompt sensitivity is evaluated using two alternatives:

\begin{itemize}
\tightlist
\item
  \texttt{An\ example\ of\ \{concept\}}
\item
  \texttt{The\ meaning\ of\ \{concept\}}
\end{itemize}

For image models, per-image embeddings are extracted and then aggregated
to concept-level representations. Image-level embeddings are cached so
that bootstrap resampling and source holdouts can be run without
recomputing forward passes.

Two layer protocols are used:

\begin{itemize}
\tightlist
\item
  \textbf{Baseline}: the selected terminal representation (\texttt{-1} /
  final-layer equivalent)
\item
  \textbf{Aligned5}: five depth fractions (\texttt{d00}, \texttt{d25},
  \texttt{d50}, \texttt{d75}, \texttt{d100})
\end{itemize}

The aligned protocol is important because it distinguishes final-output
agreement from depth-wise emergence. A convergence pattern that only
appears at terminal depth means something different from one that
accumulates monotonically across the network.

\subsubsection{3.4 Representational Similarity
Analysis}\label{representational-similarity-analysis}

RSA is computed over the 20 base concepts only.

For each model:

\begin{enumerate}
\def\labelenumi{\arabic{enumi}.}
\tightlist
\item
  Construct a 20 x 20 cosine-similarity matrix over concept embeddings.
\item
  Extract the upper triangle (190 unique concept pairs).
\item
  Compute Spearman rank correlation between model-pair vectors.
\end{enumerate}

Spearman correlation is used because the relevant question is rank
agreement in concept geometry, not absolute similarity scale, which
differs substantially across model families.

\subsubsection{3.5 Statistical Inference and
Robustness}\label{statistical-inference-and-robustness}

The final replication supplements raw RSA with four robustness
procedures.

\textbf{Mantel significance.}

\begin{itemize}
\tightlist
\item
  3,000 permutations per model pair
\item
  Two-sided p-values
\item
  Benjamini-Hochberg FDR correction across all 231 model pairs
\end{itemize}

\textbf{Image bootstrap.}

\begin{itemize}
\tightlist
\item
  300 bootstrap draws
\item
  sample size = 10 images per concept
\item
  sampling with replacement
\end{itemize}

\textbf{Source holdout.}

\begin{itemize}
\tightlist
\item
  leave-one-source-out (LOSO)
\item
  source-only analysis when enough concepts remain
\end{itemize}

\textbf{Prompt sensitivity.}

\begin{itemize}
\tightlist
\item
  three text templates for all language models
\end{itemize}

These controls are not afterthoughts. They materially shift conclusions.
The paper's main claims are based on what remains stable after these
checks.

\begin{center}\rule{0.5\linewidth}{0.5pt}\end{center}

\subsection{4. Results}\label{results}

\subsubsection{4.1 Global Geometry Is Family-Structured, Not Universally
Shared}\label{global-geometry-is-family-structured-not-universally-shared}

The most stable finding in the final replication is modality asymmetry.

\textbf{Table 2: Pairwise RSA summary by family.}

{\def\LTcaptype{none} % do not increment counter
\begin{longtable}[]{@{}
  >{\raggedright\arraybackslash}p{(\linewidth - 8\tabcolsep) * \real{0.1471}}
  >{\raggedleft\arraybackslash}p{(\linewidth - 8\tabcolsep) * \real{0.1029}}
  >{\raggedleft\arraybackslash}p{(\linewidth - 8\tabcolsep) * \real{0.2206}}
  >{\raggedleft\arraybackslash}p{(\linewidth - 8\tabcolsep) * \real{0.1912}}
  >{\raggedleft\arraybackslash}p{(\linewidth - 8\tabcolsep) * \real{0.3382}}@{}}
\toprule\noalign{}
\begin{minipage}[b]{\linewidth}\raggedright
Pair type
\end{minipage} & \begin{minipage}[b]{\linewidth}\raggedleft
Pairs
\end{minipage} & \begin{minipage}[b]{\linewidth}\raggedleft
Median \(\rho\)
\end{minipage} & \begin{minipage}[b]{\linewidth}\raggedleft
Mean \(\rho\)
\end{minipage} & \begin{minipage}[b]{\linewidth}\raggedleft
Significant after FDR
\end{minipage} \\
\midrule\noalign{}
\endhead
\bottomrule\noalign{}
\endlastfoot
Language-Language & 28 & 0.661 & 0.670 & 28 / 28 \\
Vision-Vision & 45 & 0.682 & 0.630 & 43 / 45 \\
Vision-Vision-Language & 40 & 0.702 & 0.681 & 40 / 40 \\
Language-Vision & 80 & 0.259 & 0.246 & 35 / 80 \\
Language-Vision-Language & 32 & 0.230 & 0.232 & 5 / 32 \\
Vision-Language-Vision-Language & 6 & 0.939 & 0.927 & 6 / 6 \\
\end{longtable}
}

Three features stand out.

First, within-family agreement is strong in both language and vision.
Language models agree with each other at roughly the same level as
vision models agree with each other. This supports a weak form of
convergence: once modality and objective family are held roughly fixed,
representational geometry becomes reproducible.

Second, the vision-language encoders are much closer to vision than to
language. The median vision-to-vision-language RSA (0.702) is more than
three times the median language-to-vision-language RSA (0.230). That is
the opposite of what one would expect if vision-language encoders served
as genuinely intermediate representations between the two modalities.

Third, cross-modal signal exists but remains partial. Language-vision
pairs are not centered at zero; many are significant. But they are far
below within-family agreement. The strongest language-vision pair is
I-JEPA with SmolLM3-3B-8bit (\(\rho = 0.500\)), while the strongest
vision-vision pair is ConvNeXt-v2 with Hiera-base (\(\rho = 0.929\)).
The scale of the difference matters more than the existence of a few
positive pairs.

\subsubsection{4.2 Vision-Language Models Behave Like Vision
Models}\label{vision-language-models-behave-like-vision-models}

The bridge-model question is central because multimodal models are often
used as intuitive support for Platonic convergence. The final
replication does not support that interpretation.

All 40 vision-to-vision-language pairs are significant after FDR
correction, with a median \(\rho = 0.702\). By contrast, only 5 of 32
language-to-vision-language pairs survive correction, with a median
\(\rho = 0.230\).

Representative high-similarity bridge pairs include:

\begin{itemize}
\tightlist
\item
  CLIP-ViT-B32 with I-JEPA: \(\rho = 0.823\)
\item
  I-JEPA with SigLIP2: \(\rho = 0.815\)
\item
  CLIP-ViT-B32 with Hiera-base: \(\rho = 0.776\)
\item
  SigLIP2 with ViT-MAE-base: \(\rho = 0.774\)
\end{itemize}

This pattern is better explained as engineered alignment layered on top
of a fundamentally visual encoder than as spontaneous convergence of
language and vision onto one geometry. Put differently: contrastive
training creates a useful cross-modal interface, but it does not make
the image encoder structurally language-like.

\subsubsection{4.3 Dataset Composition Changes the
Geometry}\label{dataset-composition-changes-the-geometry}

The final V2 package added explicit source holdout analysis. This was
scientifically important, because source dependence turned out to be
large enough to alter interpretation.

\textbf{Table 3: Source holdout effects relative to the full baseline
run.}

{\def\LTcaptype{none} % do not increment counter
\begin{longtable}[]{@{}
  >{\raggedright\arraybackslash}p{(\linewidth - 8\tabcolsep) * \real{0.1800}}
  >{\raggedleft\arraybackslash}p{(\linewidth - 8\tabcolsep) * \real{0.1900}}
  >{\raggedleft\arraybackslash}p{(\linewidth - 8\tabcolsep) * \real{0.2000}}
  >{\raggedleft\arraybackslash}p{(\linewidth - 8\tabcolsep) * \real{0.2200}}
  >{\raggedleft\arraybackslash}p{(\linewidth - 8\tabcolsep) * \real{0.2100}}@{}}
\toprule\noalign{}
\begin{minipage}[b]{\linewidth}\raggedright
Holdout condition
\end{minipage} & \begin{minipage}[b]{\linewidth}\raggedleft
Concepts retained
\end{minipage} & \begin{minipage}[b]{\linewidth}\raggedleft
Mean \(\Delta \rho\)
\end{minipage} & \begin{minipage}[b]{\linewidth}\raggedleft
Mean \(|\Delta \rho|\)
\end{minipage} & \begin{minipage}[b]{\linewidth}\raggedleft
Max \(|\Delta \rho|\)
\end{minipage} \\
\midrule\noalign{}
\endhead
\bottomrule\noalign{}
\endlastfoot
Leave out ImageNet & 10 & -0.1564 & 0.2358 & 0.7080 \\
Leave out OpenImages & 14 & -0.0792 & 0.0865 & 0.2654 \\
Leave out Unsplash & 16 & +0.0900 & 0.1295 & 0.4170 \\
ImageNet only & 10 & -0.0194 & 0.1912 & 0.5931 \\
\end{longtable}
}

The ImageNet result is the most consequential. Excluding ImageNet
reduces pairwise RSA by -0.1564 on average, with some model pairs
shifting by more than 0.70. This is too large to treat as nuisance
noise. It means the measured geometry is conditional on how the concept
set is instantiated in images.

This does not invalidate the experiment. It clarifies its scope. The
observed geometry is not only a property of model weights; it is a
property of model weights interacting with a concrete sampling of the
world.

\subsubsection{4.4 Prompt Choice and Uncertainty Are Not
Negligible}\label{prompt-choice-and-uncertainty-are-not-negligible}

Language-side prompt choice also matters more than a casual evaluation
might suggest.

Across the 8 language models, the mean maximum absolute shift in
cross-modal RSA under the two alternative templates is 0.1960, with a
range from 0.1446 to 0.2785. The most prompt-sensitive model in the
final run is Qwen3-0.6B-MLX-8bit; the least sensitive is
Granite-3.3-2B-Instruct-8bit.

Bootstrap uncertainty confirms that the image side contributes
substantial variance. Across the 231 model pairs, image-bootstrap
confidence intervals have:

\begin{itemize}
\tightlist
\item
  mean width: 0.1787
\item
  median width: 0.1900
\item
  maximum width: 0.3703
\end{itemize}

The implication is methodological. Small cross-modal differences should
not be overinterpreted without uncertainty estimates, because both
prompt framing and image resampling can move the results by amounts
comparable to many headline pairwise gaps.

\subsubsection{4.5 Deeper Layers Converge More, but the Final Answer
Stays the
Same}\label{deeper-layers-converge-more-but-the-final-answer-stays-the-same}

The aligned-layer analysis shows a clear depth trend:

\textbf{Table 4: Mean pairwise RSA by aligned depth fraction.}

{\def\LTcaptype{none} % do not increment counter
\begin{longtable}[]{@{}lrr@{}}
\toprule\noalign{}
Depth fraction & Mean pairwise \(\rho\) & Median pairwise \(\rho\) \\
\midrule\noalign{}
\endhead
\bottomrule\noalign{}
\endlastfoot
d00 & 0.3424 & 0.2491 \\
d25 & 0.4149 & 0.3356 \\
d50 & 0.4345 & 0.3857 \\
d75 & 0.4403 & 0.4067 \\
d100 & 0.4639 & 0.4346 \\
\end{longtable}
}

Convergence strengthens monotonically with depth. This is one of the
clearest positive findings in the entire project. It suggests that
representational agreement is accumulated, not merely read out at the
end.

But the second half of the depth result is equally important: terminal
aligned and baseline runs are almost identical. The mean
aligned-minus-baseline pairwise delta is +0.0006, and the FDR
significance profile is unchanged. So depth changes when the geometry
becomes visible, but it does not reverse the final high-level
conclusion. Deeper layers amplify the same family structure rather than
revealing a hidden universal geometry that the baseline missed.

\subsubsection{4.6 Compositionality Is Exploratory, Not a Main
Claim}\label{compositionality-is-exploratory-not-a-main-claim}

The earlier small-scale analyses encouraged a strong story about
compositionality. The final 22-model package does not support such
strong wording.

Using the current V2 metrics:

\begin{itemize}
\tightlist
\item
  mean additive compositionality is 0.948 for vision-language models
\item
  mean additive compositionality is 0.873 for vision models
\item
  mean additive compositionality is 0.831 for language models
\end{itemize}

The balance-style score is even less separating, with family means
compressed into the range 0.907 to 0.965. In other words,
compound-concept behavior in the final pipeline does not cleanly
partition by modality. It overlaps broadly and depends strongly on
metric design.

This does not make the compositionality analysis useless. It makes it
secondary. The final paper therefore treats compositionality as
exploratory evidence rather than a core pillar of the conclusion.

\begin{center}\rule{0.5\linewidth}{0.5pt}\end{center}

\subsection{5. Discussion}\label{discussion}

\subsubsection{5.1 What Converges, and What Does
Not}\label{what-converges-and-what-does-not}

The final replication supports two statements and rejects one.

Supported:

\begin{itemize}
\tightlist
\item
  models converge strongly within family
\item
  vision-language encoders converge strongly with vision encoders
\end{itemize}

Not supported:

\begin{itemize}
\tightlist
\item
  all sufficiently capable modalities converge toward one shared
  geometry
\end{itemize}

This is not a null result. It narrows the hypothesis. The data are
consistent with modality-shaped structure that remains partially coupled
across modalities but does not collapse into one shared space at the
tested scale.

\subsubsection{5.2 Why Bridge Models Do Not Rescue the Strong
Hypothesis}\label{why-bridge-models-do-not-rescue-the-strong-hypothesis}

One tempting reply is that the vision-language models already
demonstrate cross-modal convergence. But that reading confuses interface
alignment with internal geometry.

The bridge models are trained precisely to connect images and text. If
their image encoders still end up substantially closer to vision than to
language, then multimodal success cannot be taken as independent
evidence for a strong Platonic claim. It demonstrates engineered
compatibility, which is useful, but conceptually different.

\subsubsection{5.3 Why Robustness Controls Change the
Interpretation}\label{why-robustness-controls-change-the-interpretation}

The most important difference between the early prototype and the final
replication is not model count. It is methodological pressure.

Three controls matter especially:

\begin{itemize}
\tightlist
\item
  source holdout shows that image provenance materially changes geometry
\item
  prompt sensitivity shows that language-side framing materially changes
  geometry
\item
  aligned-layer analysis shows that convergence is depth-dependent but
  terminally stable
\end{itemize}

This is why the current paper is more conservative than the early
narrative. Once those controls are in place, the strongest claims are
about boundary conditions, not about universal agreement.

\subsubsection{5.4 Limitations}\label{limitations}

The paper still has clear limits.

\begin{enumerate}
\def\labelenumi{\arabic{enumi}.}
\tightlist
\item
  \textbf{No controlled retraining.} We compare pretrained models; we do
  not intervene on data or objective while holding architecture fixed.
\item
  \textbf{Scale is modest.} The language side is capped at 4B parameters
  by design. A stronger Platonic effect may emerge later.
\item
  \textbf{Concepts are still concrete-heavy.} The evaluation does not
  yet include enough abstract concepts to test whether abstraction
  strengthens or weakens cross-modal agreement.
\item
  \textbf{Prompting is narrow.} We use short definitional prompts, not
  long contextual or sensory prompts.
\item
  \textbf{Quantization sensitivity is incomplete.} A dedicated q4 vs q8
  final ablation was identified as still worth doing but is not part of
  the sign-off package.
\end{enumerate}

\subsubsection{5.5 Future Directions}\label{future-directions}

The most informative next experiments are not simply ``more models.''
They are more controlled models.

High-value next steps are:

\begin{itemize}
\tightlist
\item
  same-family scaling to isolate size from architecture
\item
  intervention on image source composition
\item
  abstract-concept expansion
\item
  richer text prompt families
\item
  controlled language-alignment training to distinguish engineering from
  emergence
\end{itemize}

If the Platonic hypothesis is true in a strong sense, it should survive
those interventions. If it does not, then the field needs a weaker and
more precise formulation.

\begin{center}\rule{0.5\linewidth}{0.5pt}\end{center}

\subsection{6. Conclusion}\label{conclusion}

We tested the Platonic Representation Hypothesis using the final local
22-model replication rather than the earlier prototype reports. The
result is not that cross-modal geometry is absent. It is that the stable
signal is weaker and more conditional than an optimistic reading
suggests.

The final evidence supports four conclusions:

\begin{enumerate}
\def\labelenumi{\arabic{enumi}.}
\tightlist
\item
  \textbf{Within-family convergence is robust.} Language models agree
  strongly with language models, and vision models agree strongly with
  vision models.
\item
  \textbf{Vision-language encoders are structurally vision-like.} Their
  geometry is much closer to vision than to language.
\item
  \textbf{Cross-modal geometry is real but partial.} It is substantially
  weaker than within-family agreement and often disappears under
  stricter significance thresholds.
\item
  \textbf{Interpretation depends on controls.} Image source, prompt
  template, and layer depth all matter enough to change the story.
\end{enumerate}

The strong Platonic claim is therefore not supported at this scale. The
more defensible statement is weaker: representational geometry is shaped
by both world structure and modality-specific training pressures, and
the balance between those forces remains an open empirical question.

\subsubsection{What We Still Do Not
Know}\label{what-we-still-do-not-know}

Two possibilities remain live.

\begin{itemize}
\tightlist
\item
  \textbf{Weak Platonic view}: the present runs are below the scale
  where cross-modal geometry becomes dominant, but the monotonic depth
  trend and non-zero cross-modal signal are early evidence.
\item
  \textbf{Family-structured view}: each modality learns a partially
  overlapping but fundamentally distinct geometry, and explicit
  alignment is required to bridge them.
\end{itemize}

The current project resolves neither possibility completely. What it
does provide is a cleaner map of the boundary: where convergence is
already robust, where it is still fragile, and which controls must be
passed before stronger claims become credible.

\begin{center}\rule{0.5\linewidth}{0.5pt}\end{center}

\subsection{References}\label{references}

Andreas, J. (2019). Measuring Compositionality in Representation
Learning. \emph{ICLR 2019}.

Bao, H., Dong, L., Piao, S., and Wei, F. (2022). BEiT: BERT Pre-Training
of Image Transformers. \emph{ICLR 2022}.

He, K., Chen, X., Xie, S., Li, Y., Dollar, P., and Girshick, R. (2022).
Masked Autoencoders Are Scalable Vision Learners. \emph{CVPR 2022}.

Huh, M., Cheung, B., Wang, T., and Isola, P. (2024). The Platonic
Representation Hypothesis. \emph{arXiv preprint arXiv:2405.07987}.

Kriegeskorte, N., Mur, M., and Bandettini, P. A. (2008).
Representational similarity analysis: connecting the branches of systems
neuroscience. \emph{Frontiers in Systems Neuroscience}, 2, 4.

Mikolov, T., Sutskever, I., Chen, K., Corrado, G. S., and Dean, J.
(2013). Distributed representations of words and phrases and their
compositionality. \emph{NeurIPS 2013}.

Oquab, M., Darcet, T., Moutakanni, T., Vo, H., Szafraniec, M., Khalidov,
V., and Bojanowski, P. (2023). DINOv2: Learning Robust Visual Features
without Supervision. \emph{arXiv preprint arXiv:2304.07193}.

Radford, A., Kim, J. W., Hallacy, C., Ramesh, A., Goh, G., Agarwal, S.,
Sastry, G., Askell, A., Mishkin, P., Clark, J., Krueger, G., and
Sutskever, I. (2021). Learning Transferable Visual Models From Natural
Language Supervision. \emph{ICML 2021}.

Zhai, X., Mustafa, B., Kolesnikov, A., and Beyer, L. (2023). Sigmoid
Loss for Language Image Pre-Training. \emph{ICCV 2023}.

\begin{center}\rule{0.5\linewidth}{0.5pt}\end{center}

\subsection{Appendix A: Statistical
Summary}\label{appendix-a-statistical-summary}

\begin{itemize}
\tightlist
\item
  Models tested: 22
\item
  Model pairs tested: 231
\item
  Base concepts used for RSA: 20
\item
  Compound concepts used for exploratory analysis: 8
\item
  Images per concept: 30
\item
  Total concept images: 840
\item
  Mantel permutations per pair: 3,000
\item
  Bootstrap draws: 300
\item
  Bootstrap image sample size: 10
\item
  Hardware target: Apple Silicon M1, 16GB RAM
\end{itemize}

\subsection{Appendix B: Reproducibility
Notes}\label{appendix-b-reproducibility-notes}

The final paper is based on the completed baseline and aligned
replication artifacts in the repository's
\texttt{experiments/01\_embeddings\_convergence\_basics/results/replication/}
outputs. The final claims in this draft are aligned with the signed-off
baseline, aligned5, and robustness artifacts rather than with the
earlier prototype summaries.

\end{document}
